\documentclass[11pt,a4paper]{article}
\usepackage[utf8]{inputenc}
\usepackage[T1]{fontenc}
\usepackage[spanish,es-noshorthands,es-nodecimaldot]{babel}
\usepackage{lmodern,microtype}
\usepackage[margin=2.2cm]{geometry}
\usepackage{xcolor,tabularx,booktabs,enumitem}
\usepackage[most]{tcolorbox}
\usepackage{listings}
\usepackage[hidelinks]{hyperref}
\hypersetup{pdftitle={Parte 4 --- ChemBERTa-2}}

\definecolor{dtured}{RGB}{153,0,0}
\definecolor{idea}{RGB}{0,90,160}
\newtcolorbox{concl}{colback=idea!6,colframe=idea,boxrule=0.6pt,left=5pt,right=5pt,top=4pt,bottom=4pt,breakable}
\newcommand{\var}[1]{\texttt{#1}}
\lstset{basicstyle=\ttfamily\footnotesize,columns=fullflexible,breaklines=true,
  literate={á}{{\'a}}1 {é}{{\'e}}1 {í}{{\'i}}1 {ó}{{\'o}}1 {ú}{{\'u}}1 {ñ}{{\~n}}1}
\setlist{itemsep=1pt,topsep=3pt}
\renewcommand{\arraystretch}{1.15}
\newcommand{\exper}[2]{\subsection*{#1}\noindent\textbf{\small\textcolor{dtured}{#2}}\par\smallskip}

\begin{document}

\begin{center}
{\LARGE\bfseries Parte 4 --- ChemBERTa-2}\\[4pt]
{\large Cambios en el código, ejecuciones realizadas y conclusiones}
\end{center}
\vspace{6pt}\hrule\vspace{10pt}

% =====================================================================
\section{Cambios respecto al código que pasaste}

Tu versión funcionaba: monta el modelo y devuelve el tamaño correcto. Los cambios son estos.

\begin{enumerate}
  \item \textbf{\var{add\_pooling\_layer=False}} al cargar el cuerpo. Esa capa no viene en los
        checkpoints, así que nacía al azar, y el \var{forward} no la usa. Con el arreglo el cuerpo se
        cuenta como 3.28\,M pesos (los reales) en vez de 3.43\,M, y desaparece el aviso \var{MISSING}.
  \item \textbf{\var{pool()} lanza error} ante un \var{pooling} desconocido. Antes un error de
        escritura (\var{CLS} en mayúsculas, por ejemplo) caía silenciosamente en la media: habríamos
        lanzado un experimento creyendo que usaba CLS cuando usaba otra cosa.
  \item \textbf{Nuevo \var{pooling: concat}}: concatena CLS y media, la cabeza recibe 768 en vez de
        384 (el tamaño de entrada se ajusta solo).
  \item \textbf{Nuevo ajuste \var{pretrained}} (defecto \var{true}). Con \var{false} monta la misma
        arquitectura con \textbf{pesos al azar}. No preentrena nada: sirve para medir cuánto aporta
        el preentrenamiento, y tarda lo mismo.
        \emph{No es lo mismo que \var{freeze\_encoder}}: \var{pretrained} fija con qué valores
        \textbf{arranca} el cuerpo, \var{freeze\_encoder} si esos valores \textbf{cambian}.
  \item \textbf{Nuevo ajuste \var{augment}} (defecto \var{false}): aleatorización de SMILES. Renumera
        los átomos al azar y escribe sin canonizar. Se aplica \textbf{solo en entrenamiento} (el
        cargador lo activa únicamente cuando \var{shuffle=True}), así validación y test siguen viendo
        el SMILES canónico.
  \item \textbf{Cambiado el checkpoint por defecto a \var{-MTR}} (ver experimento~A).
\end{enumerate}

\noindent\textbf{Archivo nuevo:} \var{batch.sh}, que ejecuta varios entrenamientos dentro de un mismo
trabajo de la cola.

\begin{concl}
\textbf{Dos cosas que descubrimos por el camino y conviene saber:}
\begin{itemize}
  \item Hay que dar a cada ejecución su \var{hydra.run.dir} con nombre. Por defecto la carpeta es la
        fecha con segundos, y con 16 trabajos a la vez \textbf{dos que arranquen en el mismo segundo
        comparten carpeta y se pisan}.
  \item La cola es el cuello de botella, no el cálculo: un día esperamos 65 minutos para 12 de
        cómputo. Por eso \var{batch.sh}, que agrupa ejecuciones en un solo trabajo.
\end{itemize}
\end{concl}

% =====================================================================
\section{Ejecuciones}

Todo en \textbf{validación}, semilla 42 salvo donde se indique, 20 épocas, cuerpo sin congelar.
Tiempos: $\approx$5 min por ejecución afinada, $\approx$1 min congelada.

% ---------------------------------------------------------------------
\exper{A. Rejilla checkpoint $\times$ pooling, más \emph{ablation} sin preentrenar (16 ejecuciones)}
{Por qué: hay que elegir checkpoint y pooling. Se cruzan en vez de fijar uno porque los dos efectos pueden estar mezclados: el CLS del MTR fue entrenado para resumir y el del MLM no (ver conclusiones), así que comparar los checkpoints con un pooling fijo mediría dos cosas a la vez. El \emph{ablation} mide cuánto aporta el preentrenamiento.}

\begin{center}
\begin{tabular}{lcc}
\toprule
Configuración & Regresión (ST-RAE$\downarrow$) & Clasificación (MCC$\uparrow$) \\ \midrule
\var{mlm\_cls}    & 0.905 & 0.142 \\
\var{mlm\_mean}   & 0.833 & 0.120 \\
\var{mlm\_concat} & 0.820 & 0.149 \\ \midrule
\var{mtr\_cls}    & \textbf{0.741} & 0.273 \\
\var{mtr\_mean}   & 0.753 & \textbf{0.292} \\
\var{mtr\_concat} & \textbf{0.741} & 0.284 \\ \midrule
sin preentrenar + \var{cls}  & 0.982 & 0.000 \\
sin preentrenar + \var{mean} & 0.817 & 0.000 \\
\bottomrule
\end{tabular}
\end{center}

\begin{concl}
\begin{itemize}
  \item \textbf{MTR gana a MLM siempre y por mucho}, en las dos tareas y con los tres \var{pooling}.
  \item \textbf{Sin preentrenar no aprende nada}: 0.982 en regresión es exactamente lo que da
        predecir siempre la media (0.983), y en clasificación 0.000 es decir ``no'' a todo.
        \textbf{Todo lo que hace ChemBERTa viene del preentrenamiento.}
  \item Dato incómodo pero honesto: \textbf{sin preentrenar + media (0.817) supera a MLM + cls
        (0.905)}. Un cuerpo al azar bate a uno preentrenado mal aprovechado.
  \item Con MTR, los tres \var{pooling} empatan (0.741 / 0.753 / 0.741).
\end{itemize}
\end{concl}

% ---------------------------------------------------------------------
\exper{B. Lo mismo con el cuerpo congelado (12 ejecuciones)}
{Por qué: congelar cuesta 5 veces menos tiempo, así que interesa saber si pierde calidad. Además, con el cuerpo quieto no hay 20 épocas para compensar un mal resumen, así que las diferencias entre pooling deberían notarse más.}

\begin{center}
\begin{tabular}{lcc@{\qquad}cc}
\toprule
& \multicolumn{2}{c}{Regresión (ST-RAE$\downarrow$)} & \multicolumn{2}{c}{Clasificación (MCC$\uparrow$)} \\
\cmidrule(r){2-3}\cmidrule(l){4-5}
Configuración & afinado & congelado & afinado & congelado \\ \midrule
\var{mlm\_cls}    & 0.905 & 0.929 & 0.142 & 0.136 \\
\var{mlm\_mean}   & 0.833 & 0.814 & 0.120 & 0.196 \\
\var{mlm\_concat} & 0.820 & 0.836 & 0.149 & 0.193 \\
\var{mtr\_cls}    & 0.741 & 0.760 & 0.273 & 0.286 \\
\var{mtr\_mean}   & 0.753 & 0.754 & 0.292 & 0.147 \\
\var{mtr\_concat} & 0.741 & 0.738 & 0.284 & 0.258 \\
\bottomrule
\end{tabular}
\end{center}

\begin{concl}
\textbf{No se puede decidir con esto.} Las diferencias en regresión van de 0.001 a 0.024, y
\textbf{la nota de validación ya oscila sola 0.026--0.039 entre épocas} de una misma ejecución
(0.027--0.070 en clasificación). Son más pequeñas que el ruido. Se repitió con semillas en E.
\end{concl}

% ---------------------------------------------------------------------
\exper{C. ¿Mejora el MCC eligiendo mejor el umbral? (sin GPU)}
{Por qué: el MCC usa un umbral fijo de 0.5. Si el problema fuera de calibración, elegir mejor el umbral debería recuperar MCC sin reentrenar nada.}

\textbf{Método:} sobre las predicciones ya guardadas. Validación se parte en dos mitades al azar: el
umbral se elige en una y el MCC se mide en la otra. 20 particiones distintas, promediadas.

\begin{center}
\begin{tabular}{lcccc}
\toprule
Ejecución & umbral 0.5 & max-MCC & Youden & prevalencia \\ \midrule
\var{mtr\_cls}        & 0.263 & 0.238 & 0.256 & 0.274 \\
\var{mtr\_mean}       & 0.284 & 0.258 & 0.239 & 0.229 \\
\var{mtr\_concat}     & 0.274 & 0.217 & 0.258 & 0.264 \\
\var{froz\_mtr\_cls}  & 0.283 & 0.239 & 0.244 & 0.245 \\
\var{froz\_mtr\_mean} & 0.158 & 0.250 & 0.275 & 0.269 \\ \midrule
\textbf{media} & 0.252 & \textbf{0.240} & \textbf{0.254} & \textbf{0.256} \\
\textbf{rango peor--mejor} & \textbf{0.126} & 0.041 & \textbf{0.036} & 0.045 \\
\bottomrule
\end{tabular}
\end{center}

\begin{concl}
\begin{itemize}
  \item \textbf{Descartado:} ninguna regla supera al 0.5 fijo de forma apreciable. Nuestros mejores
        modelos ya están bien calibrados.
  \item \textbf{Maximizar el MCC directamente es la peor opción} (0.240): prueba 97 umbrales sobre
        $\approx$115 moléculas y se ajusta al ruido. \textbf{Youden es mejor} (0.254) porque es una
        regla, no una búsqueda.
  \item \textbf{Lo interesante: Youden no sube el techo, sube el suelo.} El rango entre
        configuraciones pasa de 0.126 a 0.036. Rescata modelos mal calibrados sin tocar a los buenos.
  \item \textbf{Lectura para el informe:} el MCC con umbral fijo mezcla la capacidad de \emph{ordenar}
        con la \emph{calibración}, y penaliza más la segunda. Conviene informar MCC y AUROC juntos.
\end{itemize}
\end{concl}

% ---------------------------------------------------------------------
\exper{D. Aleatorización de SMILES y congelación, con 3 semillas (24 ejecuciones)}
{Por qué: el modelo sobreajusta desde la época 1--9 de 20, o sea que faltan datos, no capacidad. Una molécula admite muchas escrituras válidas, así que dándole una distinta cada época ve material nuevo sin medir nada en laboratorio. Es además una técnica exclusiva de los modelos de texto: Atom-JEPA no puede aplicarla. Y con 3 semillas se resuelve de paso lo que quedó sin decidir en B.}

\begin{center}
\begin{tabular}{llcc@{\qquad}cc}
\toprule
& & \multicolumn{2}{c}{Regresión} & \multicolumn{2}{c}{Clasificación} \\
\cmidrule(r){3-4}\cmidrule(l){5-6}
Cuerpo & Aleatorizar & media $\pm$ desv. & épocas pico & media $\pm$ desv. & épocas pico \\ \midrule
aprendiendo & no & \textbf{0.732 $\pm$ 0.006} & 2/2/1 & \textbf{0.292 $\pm$ 0.021} & 12/7/16 \\
aprendiendo & sí & 0.741 $\pm$ 0.005 & 1/1/1 & 0.275 $\pm$ 0.007 & \textbf{20/20/16} \\
congelado & no & 0.740 $\pm$ 0.006 & 1/1/2 & 0.277 $\pm$ 0.017 & 10/9/7 \\
congelado & sí & 0.751 $\pm$ 0.022 & 1/2/1 & 0.263 $\pm$ 0.047 & 10/6/16 \\
\bottomrule
\end{tabular}
\end{center}

\noindent Contrastes ($t$ de Student, $n=3$): aleatorizar con el cuerpo aprendiendo $p=0.13$ (reg.)
y $p=0.26$ (clas.); con el cuerpo congelado $p=0.41$ y $p=0.65$; congelar $p=0.21$ y $p=0.40$.

\begin{concl}
\begin{itemize}
  \item \textbf{Nada es significativo.} Con 3 semillas, la configuración de partida es tan buena como
        cualquier otra. \textbf{Pero las seis comparaciones apuntan al mismo lado: peor.}
  \item \textbf{Dato útil:} la dispersión \emph{entre semillas} del resultado final ($\pm0.006$ en
        regresión) es menor que la variación \emph{entre épocas} ($\pm0.026$), porque quedarse con el
        máximo de 20 es más estable que un sorteo suelto. Con 3 semillas se distinguen diferencias de
        $\approx0.01$ en regresión; en clasificación la dispersión es $\pm0.02$--$0.05$.
  \item \textbf{Cabo suelto:} con el cuerpo aprendiendo y aleatorizando, la clasificación pica en
        \textbf{20/20/16 de 20}, al borde de la ventana. Podía estar truncada. Se comprobó en E.
\end{itemize}
\end{concl}

% ---------------------------------------------------------------------
\exper{E. ¿Estaba truncada? 60 épocas en clasificación (6 ejecuciones)}
{Por qué: cerrar el cabo suelto de D. También se repite el control sin aleatorizar, porque el máximo de 60 sorteos está más inflado que el de 20 y compararlos sería injusto.}

\begin{center}
\begin{tabular}{lcc}
\toprule
Configuración & media $\pm$ desv. & época del pico \\ \midrule
sin aleatorizar, 60 épocas & \textbf{0.295 $\pm$ 0.017} & 24/7/16 de 60 \\
con aleatorizar, 60 épocas & 0.275 $\pm$ 0.007 & 20/20/26 de 60 \\
\bottomrule
\end{tabular}
\end{center}

\begin{concl}
\textbf{No estaban truncadas.} Con 40 épocas más de margen, las aleatorizadas pican en 20/20/26 y dan
\textbf{exactamente las mismas notas} que con 20 épocas (0.280 / 0.278 / 0.267). El pico en la época
20 era el techo real, no el borde de la ventana.
\textbf{La aleatorización de SMILES queda descartada}, habiendo comprobado las dos objeciones
posibles (si hacía falta que el cuerpo aprendiera, y si hacían falta más épocas).
\end{concl}

% =====================================================================
\newpage
\section{Conclusiones}

\subsection*{Lo que funciona}
\begin{enumerate}
  \item \textbf{MTR, no MLM.} Única diferencia grande y consistente: gana en las dos tareas y con los
        tres pooling, por 3--5 veces el ruido.

        \smallskip
        \textbf{Por qué, probablemente.} Inspeccionando los checkpoints, la cabeza de preentrenamiento
        del MTR es \var{dense}~$384{\to}384$ + \var{out\_proj}~$384{\to}199$: el patrón de una cabeza
        que lee \textbf{únicamente el token CLS}. Como ese era el único vector que se leía para
        predecir los 199 descriptores de RDKit, el cuerpo quedó obligado a escribir ahí un resumen de
        la molécula. En MLM la tarea leía el vector \emph{de cada posición tapada}, así que
        \textbf{el CLS no se usó nunca}. Por eso se cruzó la rejilla en A en vez de fijar un pooling.
  \item \textbf{El preentrenamiento lo es todo.} Sin él: 0.982 en regresión (= predecir la media) y
        0.000 en clasificación.
\end{enumerate}

\subsection*{Lo que se probó y se descartó, con medidas}
\begin{tabularx}{\linewidth}{@{}lX@{}}
\toprule
\var{concat} & Empata con la mejor opción simple (0.741); no justifica duplicar la cabeza (A, B) \\
Umbral que maximiza el MCC & Empeora (0.240 vs 0.252): se ajusta al ruido (C) \\
Umbral de Youden & Mejor que el anterior pero no supera al 0.5 fijo; sirve de seguro, no de mejora (C) \\
Aleatorización de SMILES & Sin efecto, en 2 tareas, congelado y afinado, 3 semillas y 60 épocas (D, E) \\
Congelar el cuerpo & Ligeramente peor sin significación. Pero cuesta \textbf{5 veces menos} (B, D) \\
\bottomrule
\end{tabularx}

\subsection*{Configuración final}
Elegida \textbf{solo con validación}:
\begin{center}
\var{checkpoint=...-MTR}\quad \var{pooling=cls}\quad \var{freeze\_encoder=false}\quad
\var{augment=false}\quad \var{epochs=20}
\end{center}
Se elige \var{cls} entre las tres que empatan por ser la más simple y la que el preentrenamiento del
MTR respalda (su cabeza leía ese mismo token),
\textbf{no porque sea mejor: no se puede distinguir}.

\subsection*{Resultados en test}
Test se calcula solo en cada ejecución, pero \textbf{no se consultó ni una vez} hasta cerrar la
configuración por validación.

\begin{center}
\begin{tabular}{lcccc}
\toprule
& \multicolumn{2}{c}{Regresión (ST-RAE$\downarrow$)} & \multicolumn{2}{c}{Clasificación (MCC$\uparrow$)} \\
\cmidrule(r){2-3}\cmidrule(l){4-5}
Modelo & validación & \textbf{test} & validación & \textbf{test} \\ \midrule
LightGBM (baseline) & 0.740 & 0.783 & 0.099 & 0.141 \\
\textbf{ChemBERTa} (3 semillas) & 0.732 & \textbf{0.778} & 0.292 & \textbf{0.183} \\
Atom-JEPA & 0.664 & \textbf{0.739} & 0.259 & \textbf{0.231} \\
\bottomrule
\end{tabular}
\end{center}

\begin{concl}
\textbf{En clasificación el orden se invierte entre validación y test.} En validación ChemBERTa
superaba a Atom-JEPA (0.292 vs 0.259); en test ocurre lo contrario (0.183 vs 0.231).

\smallskip
Caída de validación a test: LightGBM \emph{mejora} ($+0.042$), Atom-JEPA cae 0.028,
\textbf{ChemBERTa cae 0.108}.

\smallskip
\textbf{Por qué:} nuestra configuración se eligió comparando $\approx$64 ejecuciones en validación;
Atom-JEPA y LightGBM se lanzaron \emph{una vez} con sus ajustes por defecto. Cuanto más se elige
mirando una muestra, más se infla el resultado sobre ella. Afecta sobre todo al MCC, mucho más
ruidoso que el ST-RAE; en regresión no pasa (ChemBERTa cae 0.046, igual que LightGBM).

\smallskip
\textbf{No hay que presentar el 0.292 de validación como comparable al 0.259 de Atom-JEPA.}
La comparación válida es la de test.
\end{concl}

\subsection*{Tres cosas más que conviene declarar en el informe}
\begin{enumerate}
  \item \textbf{La validación es ruidosa:} oscila sola 0.026--0.039 (regresión) y 0.027--0.070
        (clasificación) entre épocas. Cualquier diferencia menor no es real.
  \item \textbf{Quedarse con la mejor de 20 épocas infla el resultado} $\approx$0.05 por azar. Afecta
        igual a todos los modelos, pero impide leer diferencias pequeñas como reales, y hace que un
        resultado de 60 épocas no sea comparable con uno de 20.
  \item \textbf{El reparto de CYP2D6 no está equilibrado:} 19.5\,\% de positivos en \var{train} frente
        a 27.4\,\% en \var{val} y 25.7\,\% en \var{test} (CYP3A4 sí: 21.3 / 20.4 / 22.2\,\%). El modelo
        calibra con una proporción y se evalúa con otra, así que predice pocos positivos y el umbral
        fijo lo penaliza. \textbf{Explica en parte por qué CYP2D6 va mal en todos los modelos.} No se
        puede corregir (hay que mantener el reparto con semilla 42), pero merece mencionarse.
\end{enumerate}

\end{document}
